\documentclass[11pt]{article}
\usepackage[localtheorems]{raeez-math-template}
\newtheorem{theorem}{Theorem}[section]
\newtheorem{proposition}[theorem]{Proposition}
\newtheorem{lemma}[theorem]{Lemma}
\newtheorem{definition}[theorem]{Definition}
\newtheorem{remark}[theorem]{Remark}
\newtheorem{corollary}[theorem]{Corollary}
\newtheorem{conjecture}[theorem]{Conjecture}
\newtheorem{construction}[theorem]{Construction}
\newtheorem{question}[theorem]{Question}
\newtheorem{hypothesis}[theorem]{Hypothesis}

\providecommand{\CY}{\mathrm{CY}}
\providecommand{\Cat}{\mathrm{Cat}}
\providecommand{\ChirAlg}{\mathrm{ChirAlg}}
\providecommand{\HH}{\mathrm{HH}}
\providecommand{\HC}{\mathrm{HC}}
\providecommand{\Ran}{\mathrm{Ran}}
\providecommand{\Fact}{\mathrm{Fact}}
\providecommand{\Obs}{\mathrm{Obs}}
\providecommand{\hCS}{\mathrm{hCS}}
\providecommand{\Rep}{\mathrm{Rep}}
\providecommand{\Perf}{\mathrm{Perf}}
\providecommand{\Coh}{\mathrm{Coh}}
\providecommand{\Ainf}{A_{\infty}}
\providecommand{\Linf}{L_{\infty}}
\providecommand{\bS}{\mathbb{S}}
\providecommand{\bP}{\mathbb{P}}
\providecommand{\bZ}{\mathbb{Z}}
\providecommand{\bC}{\mathbb{C}}
\providecommand{\bR}{\mathbb{R}}
\providecommand{\bN}{\mathbb{N}}
\providecommand{\cC}{\mathcal{C}}
\providecommand{\cO}{\mathcal{O}}
\providecommand{\cM}{\mathcal{M}}
\providecommand{\cZ}{\mathcal{Z}}
\providecommand{\cH}{\mathcal{H}}
\providecommand{\cW}{\mathcal{W}}
\providecommand{\cF}{\mathcal{F}}
\providecommand{\cL}{\mathcal{L}}
\providecommand{\cA}{\mathcal{A}}
\providecommand{\cD}{\mathcal{D}}
\providecommand{\cT}{\mathcal{T}}
\providecommand{\cE}{\mathcal{E}}
\providecommand{\fgl}{\mathfrak{gl}}
\providecommand{\fg}{\mathfrak{g}}
\providecommand{\CoHA}{\mathrm{CoHA}}
\providecommand{\Etwo}{E_2}
\providecommand{\Eone}{E_1}
\providecommand{\Ethree}{E_3}
\providecommand{\Einf}{E_{\infty}}
\providecommand{\Dolb}{\bar\partial}
\providecommand{\Ch}{\mathrm{Ch}}
\providecommand{\PV}{\mathrm{PV}}
\providecommand{\QME}{\mathrm{QME}}
\providecommand{\kch}{\kappa_{\mathrm{ch}}}
\providecommand{\kcat}{\kappa_{\mathrm{cat}}}
\providecommand{\kBKM}{\kappa_{\mathrm{BKM}}}
\providecommand{\kfib}{\kappa_{\mathrm{fiber}}}
\providecommand{\kanom}{\kappa_{\mathrm{anom}}}
\providecommand{\PhiFA}{\Phi^{\mathrm{FA}}{}}
\providecommand{\SpCh}{\mathrm{Sp}_{\mathrm{ch}}{}}
\providecommand{\Conf}{\mathrm{Conf}}
\providecommand{\CE}{\mathrm{CE}}
\providecommand{\Crit}{\mathrm{Crit}}
\providecommand{\Zchder}{Z^{\mathrm{der}}_{\mathrm{ch}}}

\title{Foundations of the CY-to-Chiral Construction $\Phi$}
\author{Raeez Lorgat}
\date{2026-04-22}

\begin{document}
\maketitle

\begin{abstract}
Fix $d \geq 1$ and a Calabi--Yau category $\cC$ with negative cyclic
Calabi--Yau class.  On the verified loci the CY-to-chiral construction is
the specialisation-indexed composite
\[
 \Phi_d^{(\Sigma_{d-1}, C)}
 =
 \SpCh_{\Sigma_{d-1}, C}\circ \PhiFA_d,
 \qquad
 \cC \xrightarrow{\;\PhiFA_d\;} \Fact^{E_d,\mathrm{hol}}(X)
 \xrightarrow{\;\SpCh_{\Sigma_{d-1}, C}\;}
 E_{n(d)}\text{-}\ChirAlg(C),
\]
with
\[
 n(d)=
 \begin{cases}
 \infty, & d=1,\\
 2, & d=2,\\
 1, & d\geq 3.
 \end{cases}
\]
The first stage is the native holomorphic factorisation algebra on the
CY$_d$ target.  The second stage is holomorphic factorisation homology
over an admissible $(d-1)$-cycle followed by restriction to a reference
curve.  For $d\geq 3$ the specialised chiral algebra is $E_1$; the
braided $E_2$ layer belongs to the Drinfeld centre, double, or module
category attached to it.  The chain-level object assignment and the
$(\infty,1)$-categorical functorial statement are separate mathematical
claims with equal force and different hypotheses.
\end{abstract}

\tableofcontents

% ================================================================
\section{The Calabi--Yau Datum}
% ================================================================

The construction reads the negative cyclic Calabi--Yau structure, the
Hochschild operations, and the framing data.  A chosen presentation of
the category enters only through witnesses used to build a chain model.

\begin{definition}[Smooth proper CY$_d$ category]
\label{def:cy-d-category}
A \emph{smooth proper CY$_d$ category} over a characteristic-zero field
$k$ is a $k$-linear pretriangulated dg or $\Ainf$ category $\cC$ with:
\begin{enumerate}[label=\textup{(H\arabic*)}]
 \item smoothness and properness: the diagonal bimodule is perfect and
 $\HH_\bullet(\cC)$ is perfect over $k$;
 \item connected unit: $\HH^0(\cC)=k$ on the connected component under
 consideration;
 \item a negative cyclic Calabi--Yau class
 $\eta_\cC\in \HC^-_d(\cC)^\vee$ whose Hochschild image gives a
 nondegenerate trace
 \[
  \mathrm{tr}_\cC\colon \HH_d(\cC)\longrightarrow k;
 \]
 \item a framed $d$-disk, equivalently $\bS^d$, compatibility datum for
 the Hochschild chains and the trace.
\end{enumerate}
For $\cC=\Perf(X)$ with $X$ a smooth projective CY$_d$ variety, the trace
is the Serre residue trace determined by a holomorphic volume form
$\Omega_X$.
\end{definition}

\begin{proposition}[Hochschild operation supplied to stage $1$]
\label{prop:hoch-operation-stage-one}
Let $\cC$ be as in Definition~\ref{def:cy-d-category}.  The Hochschild
cochain complex $\mathrm{CC}^\bullet(\cC,\cC)$ carries its usual
Gerstenhaber bracket.  After the CY$_d$ trace identifies cochains with
shifted chains, the deformation Lie operation consumed by the
factorisation-envelope construction has degree $1-d$ on the unshifted
Hochschild grading.  Equivalently,
$\HH^\bullet(\cC)[d-1]$ carries a degree-zero graded Lie operation.
\end{proposition}

\begin{proof}
Deligne's theorem gives the Hochschild cochains an $E_2$ structure and
therefore a Gerstenhaber bracket on cohomology.  The Calabi--Yau class
identifies Hochschild cochains with Hochschild chains shifted by $d$.
Transporting the bracket through this duality shifts the deformation
complex by $d-1$, which is the same statement as a degree $1-d$ bracket
before the shift.  In moduli language, Pantev--Toen--Vaquie--Vezzosi give
the derived moduli stack of objects a $(2-d)$-shifted symplectic form;
Calaque--Pantev--Toen--Vaquie--Vezzosi Koszul duality identifies the
corresponding functions with the same shifted Poisson operation.
\end{proof}

The degree $1-d$ operation is the input for Kontsevich--Tamarkin
$E_d$-formality.  At $d=1$ the elliptic curve case is commutative after
the lattice-VOA realisation.  At $d=2$ it gives the $E_2$ chiral surface
case.  At $d\geq 3$ it gives a native $E_d$ holomorphic factorisation
algebra before specialisation, while the chiral algebra on the curve is
$E_1$.

% ================================================================
\section{The Two-Stage Construction}
% ================================================================

\begin{definition}[Native holomorphic factorisation algebra]
\label{def:phi-fa}
Let $\cC$ be a CY$_d$ category and let $X$ be a smooth CY$_d$ target
realising the relevant geometric model of $\cC$.  On the verified loci,
the first stage
\[
 \PhiFA_d(\cC)\in \Fact^{E_d,\mathrm{hol}}(X)
\]
is the $E_d^{\mathrm{hol}}$ factorisation algebra assembled from:
\begin{enumerate}[label=\textup{(\roman*)}]
 \item the degree $1-d$ Hochschild deformation bracket of
 Proposition~\ref{prop:hoch-operation-stage-one};
 \item a chosen Kontsevich--Tamarkin formality or associator datum for
 the relevant little-disks operad;
 \item Costello--Gwilliam locality, converting the local $E_d$ algebra
 into a factorisation algebra on $X$;
 \item the Calabi--Yau form and Dolbeault data of $X$, refining the
 topological factorisation algebra to a holomorphic one.
\end{enumerate}
For $d\leq 2$ this is the established surface and curve construction.
For $d=3$ it is the witnessed chain-level construction after the
framing, analytic completion, and Costello--Li quantisation witnesses
are supplied.  For $d\geq 4$ the same display names a template until the
corresponding framing, holomorphic-twist, and completion witnesses are
constructed.
\end{definition}

\begin{definition}[Admissible chiral specialisation]
\label{def:sp-ch}
An \emph{admissible specialisation datum} on $X$ is a pair
$(\Sigma_{d-1},C)$, with $\Sigma_{d-1}\subset X$ a closed complex
$(d-1)$-cycle and $C\subset X$ a smooth reference curve, together with
the proper-support, Tor-independence, Beck--Chevalley, Fubini/envelope,
completion, and holomorphic descent witnesses needed for factorisation
homology.  For $\cF\in\Fact^{E_d,\mathrm{hol}}(X)$ set
\[
 \SpCh_{\Sigma_{d-1},C}(\cF)
 =
 \left(\int_{\Sigma_{d-1}}^{\mathrm{hol}}\cF\right)\bigg|_C .
\]
This is a holomorphic factorisation algebra on $C$ whose residual
operadic level is the number of complex factorisation directions
surviving the integration.
\end{definition}

\begin{theorem}[Two-stage factorisation]
\label{thm:two-stage}
On the verified locus of Definition~\ref{def:phi-fa}, and for an
admissible specialisation datum $(\Sigma_{d-1},C)$,
\[
 \Phi_d^{(\Sigma_{d-1},C)}(\cC)
 \simeq
 \SpCh_{\Sigma_{d-1},C}\bigl(\PhiFA_d(\cC)\bigr)
\]
as an $E_{n(d)}$ chiral algebra on $C$, where
\[
 n(d)=
 \begin{cases}
 \infty, & d=1,\\
 2, & d=2,\\
 1, & d\geq 3.
 \end{cases}
\]
The first stage is intrinsic after the formality and analytic witnesses
have been fixed.  The second stage is indexed by $(\Sigma_{d-1},C)$.
Two admissible specialisation data on the same CY$_d$ input may therefore
produce different chiral algebras on curves while sharing the same
native stage-$1$ object.
\end{theorem}

\begin{proof}
Stage $1$ is the three-step assembly in Definition~\ref{def:phi-fa}.
Kontsevich--Tamarkin formality supplies the $E_d$ algebraic structure
from the Hochschild operation; Costello--Gwilliam locality produces the
factorisation algebra; the holomorphic refinement uses the Dolbeault
operator and the CY form.  At $d=3$ the construction is available only
after the framing, anomaly, regulator, gauge-fixing, completion, and
sewing witnesses are fixed.

Stage $2$ is holomorphic factorisation homology over
$\Sigma_{d-1}$ followed by restriction to $C$.  Dunn--Lurie additivity
and the Francis--Gaitsgory/Ayala--Francis pushforward formalism identify
the residual little-disk direction.  At $d=1$ no transverse integration
is present and the elliptic lattice output is $E_\infty$.  At $d=2$ the
surface factorisation structure restricts to the $E_2$ chiral structure
on the curve.  At $d\geq 3$, integration over a complex $(d-1)$-cycle
leaves one complex ordered direction, hence $E_1$.
\end{proof}

\begin{corollary}[The centre carries the braided layer at $d\geq 3$]
\label{cor:e2-centre}
Let
\[
 A_\cC=\Phi_d^{(\Sigma_{d-1},C)}(\cC)
\]
with $d\geq 3$.  Then $A_\cC$ is an $E_1$ chiral algebra on $C$.  Its
braided structure is an object of the centre, double, or module layer,
for instance
\[
 \cZ\bigl(\Rep^{E_1}(A_\cC)\bigr)
 \quad\text{or}\quad
 \Rep^{E_2}\bigl(\Zchder(A_\cC)\bigr)
\]
on the verified Drinfeld-centre locus.  An $E_2$ multiplication on
$A_\cC$ itself is an additional datum and is absent from the
dimension-forced output.
\end{corollary}

% ================================================================
\section{Holomorphic Chern--Simons at Dimension Three}
% ================================================================

Holomorphic Chern--Simons gives the geometric model of
$\PhiFA_3$ on a CY$_3$ target.  It has its own typed input data and its
own quantisation hypotheses.

\begin{definition}[Geometric hCS input]
\label{def:hcs-input}
A holomorphic Chern--Simons input consists of a smooth complex threefold
$X$, a nowhere-vanishing holomorphic volume form
$\Omega_X\in H^0(X,K_X)$, and a finite type quadratic dg Lie or
$L_\infty$ algebra $(\fg,\langle-,-\rangle)$.  The field is
\[
 A\in \Omega^{0,\bullet}(X,\fg)[1],
\]
and the classical action is
\[
 S_{\hCS}(A)
 =
 \int_X \Omega_X\wedge
 \left(
  \frac12\langle A,\Dolb A\rangle
  +\frac16\langle A,[A,A]\rangle
 \right).
\]
For noncompact $X$ this is read locally or with compact support.  For
compact $X$ it becomes a quantum factorisation algebra only after the
compact package of Hypothesis~\ref{hyp:compact-hcs-package} is supplied.
\end{definition}

\begin{hypothesis}[Compact CY$_3$ hCS package]
\label{hyp:compact-hcs-package}
A compact CY$_3$ holomorphic Chern--Simons quantisation requires:
\begin{enumerate}[label=\textup{(\roman*)}]
 \item an orientation and determinant-line framing compatible with the
 BV pairing;
 \item a quadratic gauge dg Lie algebra and, when colour traces are
 used, a representation $\rho\colon\fg\to\mathrm{End}(W)$;
 \item elliptic Costello gauge fixing, heat-kernel regularisation,
 propagators, and renormalisation counterterms;
 \item completed local functionals and observables on which the BV
 Laplacian, bracket, and RG flow are defined;
 \item a solution of the scale-dependent quantum master equation;
 \item vanishing or explicit trivialisation of the anomaly class
 $\kanom(X,\fg,\rho)$, possibly after open--closed counterterms;
 \item factorisation descent over $X$ and TCFT sewing data whenever
 global traces or partition functions are asserted.
\end{enumerate}
\end{hypothesis}

\begin{proposition}[Local hCS and Hochschild operations]
\label{prop:local-hcs-hoch}
Let $U\subset\bC^3$ be a Stein polydisc with
$\Omega=dz_1\wedge dz_2\wedge dz_3$.  After a chain formality datum, the
standard Costello--Li local anomaly-free BV quantisation, and completed
local observables have been fixed, there is an $E_3^{\mathrm{hol}}$
quasi-isomorphism between local quantum hCS observables with matrix
gauge algebra $\fgl(V)$ and the matrix-framed Hochschild operations of
$\Perf(U)$:
\[
 \Obs_{\hCS}^q(U,\fgl(V))
 \simeq
 C^\bullet_{\mathrm{Hoch},E_3}\bigl(\Perf(U);\mathrm{End}(V)\bigr).
\]
At the symbol level, HKR gives
\[
 \HH^\bullet(\Perf(U))
 \simeq
 \PV^\bullet(U)
 =
 \bigoplus_{p,q}\Omega^{0,q}(U,\wedge^pT^{1,0}U),
\]
and contraction with $\Omega$ sends $\PV^{p,q}(U)$ to
$\Omega^{3-p,q}(U)$.
\end{proposition}

\begin{proof}
The HKR map identifies Hochschild cochains with polyvectors on the
affine formal model.  The CY form converts polyvectors to complementary
Dolbeault forms.  The chosen matrix object supplies the open sector
$\mathrm{End}(V)$ and therefore the hCS dg Lie algebra
$\Omega^{0,\bullet}(U,\fgl(V))[1]$.  Costello--Gwilliam locality turns
the anomaly-free BV package into a factorisation algebra.  The
open--closed map identifies closed Hochschild operations with local
operations on this hCS sector.
\end{proof}

\begin{proposition}[Bochner--Martinelli two-point kernel]
\label{prop:bm-kernel}
On the flat abelian model $U=\bC^3$, the propagator entering the
one-loop OPE is represented by the Bochner--Martinelli kernel
\[
 P_{\mathrm{BM}}(z,w)
 =
 \frac{2}{(2\pi i)^3}
 \sum_{k=1}^{3}(-1)^{k-1}
 \frac{\overline{(z_k-w_k)}}{\|z-w\|^6}
 \widehat{d\bar z_k}\wedge dw_1\wedge dw_2\wedge dw_3 .
\]
The corresponding observable algebra is scheme-independent only after
the QME, counterterm, and gluing package of
Hypothesis~\ref{hyp:compact-hcs-package} has selected a completed
quantum model.
\end{proposition}

\begin{definition}[Open hCS anomaly slot]
\label{def:hcs-anomaly-slot}
For a representation $\rho\colon\fg\to\mathrm{End}(W)$ define
\[
 I_4^\rho(x_1,x_2,x_3,x_4)
 =
 \frac1{4!}\sum_{\sigma\in S_4}
 \mathrm{Tr}_W\!\left(
  \rho(x_{\sigma(1)})\rho(x_{\sigma(2)})
  \rho(x_{\sigma(3)})\rho(x_{\sigma(4)})
 \right).
\]
The nonabelian one-loop gauge contribution to the hCS QME in complex
dimension three is locally represented by
\[
 \Theta^\rho_U(\alpha_0,\alpha_1,\alpha_2,\alpha_3)
 =
 \int_U
 \left[
 I_4^\rho
  \bigl(\alpha_0,\partial\alpha_1,\partial\alpha_2,\partial\alpha_3\bigr)
 \right]_{(3,3)}.
\]
Its class
\[
 \kanom(X,\fg,\rho)
 =
 [\,\hbar\,\Theta^\rho\,]
 \in
 H^1_{\mathrm{loc}}
 \bigl(\Omega^{0,\bullet}(X,\fg)[1],\cO_{\mathrm{loc}}\bigr)
\]
is the first open-string obstruction to the QME.
\end{definition}

The quadratic Casimir controls field-strength renormalisation.  The
quartic invariant $I_4^\rho$ controls the degree-one QME obstruction.
These invariants sit in different graph arities and different
cohomological degrees.

% ================================================================
\section{The Computational Chain-Level Reading}
% ================================================================

The one-dimensional calculation attached to a specialisation datum is
the chain model of the second stage.

\begin{construction}[Cyclic-to-chiral chain]
\label{constr:cyclic-to-chiral}
After a formality datum and an admissible specialisation datum have been
fixed, the chain-level construction proceeds as:
\[
 \text{cyclic }\Ainf\text{ input}
 \longrightarrow
 \text{local Lie}^{*}\text{ algebra on }C
 \longrightarrow
 \text{chiral factorisation envelope}
 \longrightarrow
 E_{n(d)}\text{ enhancement}
 \longrightarrow
 \text{quantum chiral algebra}.
\]
The negative cyclic class supplies the trace and Connes-closedness.  The
Hochschild operation supplies the local Lie structure.  The
factorisation envelope gives the chiral algebra on $\Ran(C)$.  The
framing and specialisation datum set the operadic level.  The quantum
deformation is controlled by the BV or Hochschild obstruction complex
appropriate to the chosen model.
\end{construction}

\begin{proposition}[What the construction proves]
\label{prop:what-construction-proves}
Construction~\ref{constr:cyclic-to-chiral} proves an object-level
assignment on the loci where its witnesses are present:
\[
 \cC\longmapsto
 A_\cC^{(\Sigma_{d-1},C)}
 =
 \SpCh_{\Sigma_{d-1},C}\bigl(\PhiFA_d(\cC)\bigr).
\]
At $d=3$ this includes the chain-level $\bS^3$ framing, the compact hCS
package when a compact hCS model is used, and the filtered Hochschild
comparison that kills the primary derived obstruction.  Smooth
properness of $\cC$ alone supplies none of these witnesses.
\end{proposition}

% ================================================================
\section{Object Assignment and Functoriality}
% ================================================================

\begin{theorem}[Two scopes of the symbol $\Phi_d$]
\label{thm:two-scopes}
The notation $\Phi_d$ has two precise readings.
\begin{enumerate}[label=\textup{(\alph*)}]
 \item \emph{Object-level chain construction.}  For a fixed
 $(\Sigma_{d-1},C)$ and a supplied set of witnesses, it assigns to
 $\cC$ the concrete chiral algebra
 $A_\cC^{(\Sigma_{d-1},C)}$ of
 Proposition~\ref{prop:what-construction-proves}.
 \item \emph{$(\infty,1)$-categorical functoriality.}  It sends cyclic
 $\Ainf$ or dg functors compatible with the CY structure and framing to
 morphisms of the corresponding chiral algebras.
\end{enumerate}
The first statement is proved on the verified object-level loci.  The
second statement is a per-$d$ functoriality theorem to be proved; in
general it remains conjectural.  Neither reading is subordinate to the
other.  The object-level statement carries explicit complexes, OPE
kernels, Borcherds products, and Hodge supertraces.  The
$(\infty,1)$ statement carries morphism preservation, descent, and
centre functoriality.
\end{theorem}

\begin{conjecture}[Morphism preservation]
\label{conj:morphism-preservation}
For every cyclic $\Ainf$ functor
$f\colon\cC\to\cC'$ between smooth proper CY$_d$ categories compatible
with the chosen framings and specialisation data, the construction
induces a morphism
\[
 \Phi_d^{(\Sigma,C)}(f)\colon
 \Phi_d^{(\Sigma,C)}(\cC)\longrightarrow
 \Phi_d^{(\Sigma,C)}(\cC')
\]
functorial up to coherent homotopy.  The $d=2$ Mukai-transform case is
the first test; compact CY$_3$ cases require the same witnessed
framing, completion, and pushforward data as the object-level
construction.
\end{conjecture}

% ================================================================
\section{Five Objects Kept Separate}
% ================================================================

\begin{definition}[The five associated objects]
\label{def:five-objects}
Let $A$ be an $E_1$ chiral algebra on a curve.  The five objects
\[
 A,\qquad B(A),\qquad A^i,\qquad A^!,\qquad \Zchder(A)
\]
have distinct meanings:
\begin{enumerate}[label=\textup{(\roman*)}]
 \item $A$ is the chiral algebra itself;
 \item $B(A)$ is its bar coalgebra or chiral bar construction;
 \item $A^i$ is the Koszul dual coalgebra on the conilpotent or
 completed locus;
 \item $A^!$ is the Verdier or Koszul dual algebra, according to the
 duality context;
 \item $\Zchder(A)$ is the derived chiral centre, the bulk object whose
 representation theory carries the braided layer at $d\geq 3$.
\end{enumerate}
\end{definition}

\begin{proposition}[Bar-cobar inversion and duality]
\label{prop:bar-cobar-five}
Under the usual augmented, conilpotent, or completed hypotheses, the
bar-cobar counit gives
\[
 \Omega B(A)\simeq A .
\]
This is inversion for the bar-cobar adjunction.  The passage from $A$ to
$A^!$ uses Verdier or Koszul duality data.  The passage from $A$ to
$\Zchder(A)$ is Hochschild or centre formation.  These three operations
have different source categories and different universal properties.
\end{proposition}

% ================================================================
\section{Compact and Noncompact Domains}
% ================================================================

The compact smooth proper construction and the noncompact equivariant
construction use the same CY-to-chiral grammar with different input
hypotheses.

\begin{definition}[Noncompact equivariant CY datum]
\label{def:noncompact-cy}
Let $X$ be a smooth noncompact CY$_d$ variety with a torus action
$T$ preserving the holomorphic volume form.  The equivariant input is
\[
 \cC=\Perf^T_{\mathrm{cpt}}(X),
\]
the category of $T$-equivariant perfect complexes with compact support.
The trace is read on the fixed or compactly supported sector by
equivariant Serre duality.  For a quiver-with-potential model the
corresponding Hall algebra is computed on critical cohomology of
$\Crit(W)$.
\end{definition}

\begin{proposition}[Toric CY$_3$ positive half]
\label{prop:c3-positive-half}
For the toric CY$_3$ input $X=\bC^3$ with equivariant parameters
$\epsilon_1+\epsilon_2+\epsilon_3=0$,
\[
 \CoHA(\bC^3)\simeq Y^+(\widehat{\fgl}_1).
\]
The full $\cW_{1+\infty}$ algebra at $c=1$ is recovered after the
Drinfeld double or derived-centre passage.  Thus the direct
object-level $\Phi_3$ output on the toric positive Hall side is the
positive half $Y^+$.
\end{proposition}

\begin{remark}[Conifold and local surfaces]
\label{rem:noncompact-examples}
The resolved conifold, local $\bP^2$, and $A_n$-McKay local surfaces
belong to the noncompact equivariant input class when a compact-support
or critical-cohomology model is fixed.  The conifold has
$\kch=1$ by the direct McKay calculation.  Local $\bP^2$ has
$\kch=3/2$.  These are $\kch$ values attached to the chosen
noncompact equivariant construction.
\end{remark}

% ================================================================
\section{\texorpdfstring{$K3\times E$}{K3 x E} and Borcherds Weights}
% ================================================================

\begin{proposition}[Construction-specific invariants on $K3\times E$]
\label{prop:k3e-invariants}
For $X=K3\times E$,
\[
 \kcat(X)=\chi(\cO_{K3})\chi(\cO_E)=2\cdot0=0,
\]
and the compact Hodge supertrace gives
\[
 \kch^{\mathrm{compact}}(X)
 =
 \sum_{q=0}^{3}(-1)^qh^{0,q}(X)
 =
 1-1+1-1
 =
 0.
\]
The Heisenberg--Mukai specialisation has
\[
 \kch^{\mathrm{Heis}}(X)=3,
\]
the Gritsenko $\Delta_5$ Borcherds construction has
\[
 \kBKM(\mathfrak g_{\Delta_5})=5,
\]
and the K3 fibre Mukai rank gives
\[
 \kfib(X)=24.
\]
The four construction-specific values are therefore
\[
 \{\kcat,\kch^{\mathrm{Heis}},\kBKM,\kfib\}(K3\times E)
 =
 \{0,3,5,24\}.
\]
Each number comes from its own construction: total-space Euler
characteristic, Heisenberg--Mukai chiral rank, Borcherds weight, and
fibre Mukai lattice rank.
\end{proposition}

\begin{theorem}[Borcherds weight identity]
\label{thm:borcherds-weight}
Let $\Phi_N$ be a Borcherds product obtained from a weak Jacobi form of
weight $0$ and index $1$, with constant Fourier coefficient $c_N(0)$ in
the normalisation of the lift.  The BKM weight invariant is
\[
 \kBKM(\Phi_N)=\frac{c_N(0)}{2}.
\]
In the Gritsenko--Clery eight-row diagonal-divisor catalogue used here,
\[
 (c_i(0),\kBKM_i)_{i=1}^{8}
 =
 (10,5),(4,2),(6,3),(2,1),(4,2),(1,\tfrac12),(3,\tfrac32),(2,1).
\]
The $N=1$ K3 row gives $c_1(0)=10$ and
$\kBKM(\Delta_5)=5$.
\end{theorem}

\begin{proof}
Borcherds' singular theta lift computes the weight of the automorphic
product as one half of the constant coefficient of the input weak
Jacobi form.  Gritsenko--Nikulin apply this to the K3 elliptic genus and
obtain the $\Delta_5$ row.  The CHL and diagonal-divisor variants use
the same weight formula after inserting the corresponding twisted Jacobi
form and multiplier system.
\end{proof}

The additive expression
\[
 \kBKM \stackrel{?}{=} \kch+\chi(\cO_{\mathrm{fiber}})
\]
is unsupported by the proved identities.  The $N=1$ row already decides
the matter:
\[
 \kBKM(\Delta_5)=5,\qquad
 \kch^{\mathrm{compact}}(K3\times E)+\chi(\cO_E)=0+0=0.
\]
The Borcherds invariant is the automorphic weight $c_N(0)/2$.

% ================================================================
\section{Proof Obligations}
% ================================================================

\begin{question}[Full morphism functoriality]
\label{q:morphism-functoriality}
Does Conjecture~\ref{conj:morphism-preservation} hold for cyclic
$\Ainf$ morphisms in every verified dimension?  The Mukai transform at
$d=2$ is the first chain-level test.
\end{question}

\begin{question}[Compact CY$_3$ outside witnessed loci]
\label{q:compact-cy3-outside}
Given a compact CY$_3$ without a supplied K3-fibre specialisation,
which connective Hochschild or hCS model computes the framing
obstruction and the QME anomaly class?
\end{question}

\begin{question}[Noncompact nonequivariant CY$_3$]
\label{q:noncompact-nonequivariant}
For a smooth noncompact CY$_3$ without a torus action, can Joyce--Song
or motivic critical-cohomology methods supply the compact-support trace
and factorisation homology witnesses needed for $\Phi_3$?
\end{question}

\begin{question}[Global quantum group]
\label{q:global-quantum-group}
For which compact CY$_3$ inputs does the positive Hall side admit a
completed Drinfeld double compatible with descent, giving a construction
of the global quantum group $G(X)$?
\end{question}

% ================================================================
\section{Primary Anchors}
% ================================================================

\begin{itemize}[leftmargin=1.4em]
\item Beilinson--Drinfeld, \emph{Chiral Algebras}, 2004: chiral
factorisation envelopes and descent on curves.
\item Borcherds, automorphic products on $\mathrm{O}(2,n)$, 1995 and
1998: weight equals one half of the constant coefficient.
\item Costello, \emph{Renormalization and Effective Field Theory}, 2011:
BV quantisation and renormalisation.
\item Costello--Gwilliam, \emph{Factorization Algebras in Quantum Field
Theory}, Vols. I--II: classical and quantum factorisation algebras.
\item Costello--Li, holomorphic Chern--Simons and BCOV quantisation in
complex dimension three.
\item Francis--Gaitsgory and Ayala--Francis: factorisation homology and
pushforward/excision.
\item Gritsenko--Nikulin and Gritsenko--Clery: K3 Borcherds products,
diagonal divisors, and the $\Delta_5$ row.
\item Kontsevich--Tamarkin and Fresse--Willwacher: formality of little
disks and graph-complex control of formality choices.
\item Pantev--Toen--Vaquie--Vezzosi and Calaque--Pantev--Toen--Vaquie--
Vezzosi: shifted symplectic and shifted Poisson structures.
\item Schiffmann--Vasserot and Tsymbaliuk: $\CoHA(\bC^3)$ and the
positive half of the affine Yangian.
\end{itemize}

\end{document}
